% !TEX root = main.tex
% sets up the referencing engine with a few important things. I personally prefere to have article titles and a doi in my references for ease of perousal, but this should be discussed with your supervisor. You can easily remove those sections from the command.

% For Tetrahedron, Analytical Chemistry, or any other ACS hournal - chem-acs
% For Angewandte Chemie (and Wiley journals) - chem-angew
% For any RSC journals - chem-rsc
% For Biochemistry (and some ACS journals) - chem-biochem

% Check with the journal itself as to any referencing specifics that you'll need, mor einfo about the bibtex style customisations for chemistry can be found here: https://mirror.aarnet.edu.au/pub/CTAN/macros/latex/contrib/biblatex-contrib/biblatex-chem/biblatex-chem.pdf 

\usepackage[style=chem-acs, articletitle=true, doi=true, backend=biber]{biblatex}
\usepackage{
	url,
	csquotes
}
% add in your bibliography file
\addbibresource{/Users/adrea/library.bib}

% roman numerals for footnotes
\renewcommand{\thefootnote}{\Roman{footnote}}

% superscript referencing
\renewcommand{\cite}{\autocite}
\newcommand{\citea}[1]{\citeauthor{#1}\cite{#1}}

